\begin{myChapterIntro}{本章涵盖}
\begin{itemize}
\item 软件设计的基础
\item 良好软件设计的益处
\item 如何分析应用的需求，以设计\emph{正确的应用}
\item 如何运用良好的设计技术，\emph{把应用做对}
\end{itemize}
\end{myChapterIntro}

设计良好的程序能够完成它们应该完成的工作。它们更可靠、更灵活，也更易维护。此外，它们更易于测试，而且往往比设计糟糕的程序更早完成。设计良好的程序在很多方面都更胜一筹。

为了提升你的软件设计技能，本书将教你一系列原则和模式，使你能够开发出设计良好的可持续应用。可持续的应用（sustainable application）是寿命很长的应用，因此我们希望它可靠、灵活且易于维护。在当今竞争激烈的就业市场中，一流的设计技能受到雇主的高度追捧。你的职业生涯要求你了解并运用良好的软件设计技术。

\myGraphic{0.897}{images/CH01_00_UN01_Mak.png}{}

本书将教你面向对象的\emph{设计原则}和\emph{设计模式}，以此提升你的软件设计技能。设计原则有助于改进几行代码、一个函数、一整个类，或一组协同工作的类的设计。设计模式提供了解决常见软件架构问题的模型。这些模型以设计原则为基础。

实现设计良好的软件通常并非一帆风顺。我们首先需要获取应用的需求。为实现良好的设计，往往需要多轮开发迭代，可能还要进行一些回溯，以从糟糕的设计决策中恢复过来。设计良好的应用源自艰辛的努力。

为了从本书中获得最大收益，你至少应是初级到中级的 C++ 程序员。你应熟悉基本数据结构及其算法，并理解面向对象编程（object-oriented programming，OOP）。此外，你的编程能力应足以编写由多个头文件（.h 或 .hpp）和实现文件（.cpp）源文件组成的应用，并能够编辑、编译、调试和运行它们。

由于本书中用于良好设计的技术建立在 OOP 概念之上，我们将在本章末尾简要回顾这些概念。

\section{什么是软件设计？}

设计是一种有纪律的工程方法，用于为某个问题创建解决方案。对软件开发人员来说，解决方案就是一个满足其需求的成功应用。我们通过运用本书所介绍的设计技术来实践有纪律的软件工程，以找到从需求通往设计良好、可持续的应用的最佳解决路径。这些技术包括改进我们代码的设计原则，以及帮助解决常见软件架构问题的设计模式。

\myGraphic{0.761}{images/CH01_00_UN02_Mak.png}{}

良好的设计原则有助于消除那些糟糕的意外，比如代码行为不符合预期或性能很差。这些原则有助于让我们的代码更灵活，能够应对诸如新需求之类的变化。设计模式在更高的设计层次上运作，并建立在设计原则之上。这些模式是经过业界验证的模型，用于为常见软件架构问题创建定制解决方案。

\section{你将能从本书中学到什么}

本书面向希望学习良好软件设计技能的初级到中级软件开发人员。通过许多“良好设计之前”和“良好设计之后”的程序示例，本书将教你如何

\begin{itemize}
\item 运用设计原则来改进你的代码
\item 运用设计模式——它们是经过业界验证的模型，用于解决常见软件架构问题
\item 收集、确认并分析应用的需求，以确保你编写的是正确的应用，并将其设计好
\item 以迭代方式开发设计良好的应用，并通过回溯从糟糕的设计决策中恢复
\end{itemize}

\begin{mySidebar}{通过示例学习}

通过示例学习效果最好。如果只是有人告诉我们要使用某条设计原则或某个设计模式，效果往往并不理想。我们想看到这条原则或模式\emph{如何}以及\emph{为何}能让程序变得更好。我们将努力用一个又一个示例应用来证明每一条原则或模式。用“之前”和“之后”的程序对比，可以突出设计上的改进。
\end{mySidebar}

本书的程序示例使用 C++，一种非常流行的应用开发语言。C++ 是一门非常复杂、特性繁多的语言。为了让我们能够专注于设计原则和设计模式，程序示例只使用该语言的基本特性。

\begin{myWarning}{注意}
本书中的示例程序使用 2017 版的 C++。
\end{myWarning}

\myGraphic{0.783}{images/CH01_00_UN03_Mak.png}{}

\section{良好软件设计的益处}

可持续的应用有其生命周期，在此期间它被部署、被客户成功使用，并持续得到维护。随着应用需求的演变，新版本会修复缺陷并添加新功能。良好的软件设计能够可靠地创造出可持续的应用。

\begin{myWarning}{注意}
尽管一个应用可能需要多个协同工作的独立程序，但本书中的每个示例应用都相对简短、简单，以便于看清它所示例的技术。因此，我们常常会互换使用\emph{应用}和\emph{程序}这两个词（后者指实现该应用的程序）。
\end{myWarning}

我们想要超越那些临时凑合的写法，以及那种“不管怎样先按时交差”的编程风格。当然，只要合适，临时凑合一下并无不妥。有时你只是需要让计算机完成某件事，如果一两行代码或一小段脚本就能搞定，那就放手去做吧！

没有人会反对按时完成应用。通过把应用设计好，我们往往能更快地完成它。设计良好的应用以系统化的方式逐步演进。它能更早达到最小可行产品（minimum viable product，MVP）的状态，即通过测试并满足一整套最低限度的需求。如有必要，MVP 可以作为应用的第一个版本部署。设计良好的应用更易于测试，在你需要添加功能或做其他改动时也更灵活。

我们的目标是创造出设计良好的应用：它们满足自身需求、按时完成，并且易于维护。设计良好的软件

\begin{itemize}
\item \emph{满足自身需求}——它能完成自己应该完成的工作。
\item \emph{可靠}——它通过测试，且缺陷更少。
\item \emph{符合用户的预期}——当我们调用一个函数或创建一个对象时，不应为其结果感到意外。
\item \emph{高效}——不存在隐藏的运行时性能问题。
\item \emph{灵活且可伸缩}——当需求发生变化时，很容易添加新功能，而不会增加软件的复杂度。
\item \emph{便于协作}——开发者能够更好地协同工作，并更快地从糟糕的设计决策中恢复。
\item \emph{易于维护}——设计良好的代码更容易被后来的开发者理解。
\item \emph{运用良好的设计技术}——良好的设计技术能够改进软件，例如简化代码、消除重复代码。
\item \emph{采用恰当的设计模式}——恰当的设计模式是经过业界验证的模型，用于解决常见软件架构问题。
\item \emph{总体上节省时间和成本}——良好的设计能减少开发过程中的错误和重大返工。为可持续的应用做良好设计而在前期额外花费的时间，都将由部署后维护时间和成本的降低，以及可能更长的生命周期来补偿。
\item \emph{是更好的代码}——你可以为自己写出的优秀代码而自豪。
\end{itemize}

\myGraphic{0.779}{images/CH01_00_UN04_Mak.png}{}

\section{几个设计示例}

软件设计要处理许多问题以改进我们的代码。以下是几个说明性的示例。本书还会涵盖更多设计问题。

\subsection{1.4.1 变化的泄漏}

所有程序员的祸根，就是修改程序的某一部分之后，又不得不去修改其他部分。这些修改可能会级联开来，导致重写程序的大部分甚至全部。下面的代码清单是一个示例 \texttt{Car} 类，它正体现了这一问题。

\begin{codeboxt}{代码清单 1.1（程序 1.1 Changes）：\texttt{Car.h}（设计糟糕）}
class Car
{
public:
    void step_on_brake();
    void insert_key();
    void turn_key();
    void step_on_accelerator();
};
\end{codeboxt}

在代码清单 1.2 中，类 \texttt{Driver} 使用了类 \texttt{Car}。

\begin{codeboxt}{代码清单 1.2（程序 1.1 Changes）：\texttt{Driver.h}（设计糟糕）}
#include "Car.h"

class Driver
{
public:
    Driver(Car& c) : car(c) {}
    void start_car()
    {
        car.step_on_brake();
        car.insert_key();
        car.turn_key();
        car.step_on_accelerator();
    }
private:
    Car car;
};
\end{codeboxt}

如果以后我们想要更现代的汽车，就会把

\begin{codebox}
    void insert_key();
    void turn_key();
\end{codebox}

替换为

\begin{codebox}
    void press_start_button()
\end{codebox}

在类 \texttt{Car} 中，我们会被迫对类 \texttt{Driver} 做出修改。我们在类 \texttt{Car} 中所做的修改泄漏到了类 \texttt{Driver} 中。正如我们将在本书通篇的示例中看到的那样，良好的软件设计通过减少类之间的依赖，有助于防止这类泄漏。

\subsection{1.4.2 过于复杂的代码}

类 \texttt{Automobile} 很复杂，因为它试图做太多的事情。

\begin{codeboxt}{代码清单 1.3（程序 1.2 Changes）：\texttt{Automobile.h}（设计糟糕）}
class AutomobileApp
{
public:
    void accelerate();
    void adjust_headlights();
    void apply_brakes();
    void change_oil();
    void change_tires();
    void check_brakes();
    void check_tires();
    void rotate_tires();
    void shut_off_engine();
    void start_engine();
    void tuneup_engine();
    void turn_left();
    void turn_right();
    void vacuum_car();
    void wash_car();
    void wax_car();
 
private:
    Brake brakes[4];
    Engine engine;
    Oil engine_oil;
    Direction heading;
    Headlight headlights[4];
    int speed;
    Soap soap;
    Tire tires[4];
    Vacuum vacuum_cleaner;
    CarWax wax;
};
\end{codeboxt}

本书通篇的示例展示了如何通过把每个类设计为只承担一项主要职责，来避免这个非常常见的问题。

类的泛滥是程序变得过于复杂的另一种方式。图 1.1 说明了这种情况多么容易失控。

\myGraphic{0.791}{images/CH01_F01_Mak.png}{图 1.1 类与子类的层次结构。它非得这么复杂吗？良好的设计技术可以简化这些数据。}

我们真的需要所有这些类吗？良好的设计有助于从程序中消除不必要的类。不过，有时运用良好的设计技术意味着要添加类，以使程序更灵活。这两种情况我们都会看到示例。

\myGraphic{0.703}{images/CH01_F01_UN05_Mak.png}{}

\subsection{1.4.3 不灵活的代码}

假设有一个应用，其中类 \texttt{Pet} 有两个子类 \texttt{Cat} 和 \texttt{Dog}。

\begin{codeboxt}{代码清单 1.4（程序 1.3 Inflexible）：\texttt{Pet.h}}
class Pet
{
    //...
};

class Cat : public Pet
{
    //...
};

class Dog : public Pet
{
    //...
};
\end{codeboxt}

此外，假设在应用的另一个部分，有这样一条语句

\begin{codebox}
    Cat *pet = new Cat();
\end{codebox}

这条语句不灵活。如果以后我们需要让变量 \texttt{pet} 指向一个宠物 \texttt{Dog} 对象，该怎么办？或者如果我们添加了 \texttt{Hamster} 和 \texttt{Goldfish} 子类，并希望引用这些类型的宠物，又该怎么办？我们将不得不修改包含前面那条语句的源代码，因为我们在编写这条语句时就决定只引用一只猫。这被称为\emph{硬编码}——我们在编写代码时就冻结了决策。

我们将看到良好的设计技术如何鼓励我们编写灵活的代码，使决策（例如引用哪种类型的 \texttt{Pet} 对象）能够在运行时动态做出。

\subsection{1.4.4 意外！}

如果编写设计糟糕的代码，这些代码可能会给使用它的其他程序员埋下讨厌的意外，尤其是当意外便是错误结果的时候。下面这个 \texttt{Date} 类就是一个例子。

\begin{codeboxt}{代码清单 1.5（程序 1.4 Surprise）：\texttt{Date.h}（设计糟糕）}
#include <string>

using namespace std;

class Date
{
public:
    Date(int y, int m, int d) : year(y), month(m), day(d) {}
        
    string date_string() const
    {
        return MONTH_NAMES[month] + " "  + to_string(day)
                                  + ", " + to_string(year);
    }
        
private:
    const string MONTH_NAMES[12] =
    {
        "JAN", "FEB", "MAR", "APR", "MAY", "JUN",
        "JUL", "AUG", "SEP", "OCT", "NOV", "DEC"
    };
        
    int year, month, day;
};
\end{codeboxt}

下面这个代码清单中的测试程序创建并打印一个日期字符串。

\begin{codeboxt}{代码清单 1.6（程序 1.4 Surprise）：\texttt{tester.cpp}}
#include <iostream>
#include "Date.h"

using namespace std;

int main()
{
    Date date(2023, 9, 2);  //SEP 2, 2023
    cout << date.date_string() << endl;
 
    return 0;
}
\end{codeboxt}

这个程序打印出的日期字符串是什么？

\myGraphic{0.640}{images/CH01_F01_UN06_Mak.png}{}

正如本书中许多示例程序所演示的那样，设计良好的程序能够完成它应该完成的工作，不会有意外的表现。

\subsection{1.4.5 常见的架构问题}

程序员经常会遇到一些其实相当常见的软件架构问题。设计模式为其中的许多问题提供了开发定制解决方案的模型。例如，发布—订阅问题涉及一个应用组件——称为发布者（publisher），它产生数据。其他应用组件——称为订阅者（subscriber），消费这些数据（图 1.2）。我们将在本书后面看到，观察者设计模式（Observer Design Pattern）如何为我们提供模型，来为这一架构问题开发解决方案。

\myGraphic{0.460}{images/CH01_F02_Mak.png}{图 1.2 发布—订阅是一个常见的软件架构问题。一个应用组件产生数据，其他应用组件消费这些数据。观察者设计模式提供了解决这一问题的模型。}

\section{确保我们要构建的是正确的应用，然后把它构建好}

设计应用的关键起点是获取并分析它的需求，以确保我们开发的是正确的应用。一个应用无论设计得多么好，如果它无法完成自己应该完成的工作，就不算成功。获取良好的需求是第 3 章的重要主题。从需求出发，我们可以确定初始的类集合。之后我们运用良好的设计技术，把应用构建好。

\section{良好的设计来之不易}

要持续开发出设计良好的软件，需要实践和经验。开发一个应用往往需要多轮“设计—编码—测试”迭代。我们必须做出设计上的权衡，并且可能需要从糟糕的设计决策中回溯。通往良好设计的道路崎岖不平。如果我们不是迭代式地开发软件，而是试图用一场旷日持久的编码马拉松把它一次性完成，那就不该指望一个能正常工作的应用会在结束时凭空神奇地出现。这种“大爆炸”几乎从不发生（图 1.3）。

\myGraphic{0.753}{images/CH01_F03_Mak.png}{图 1.3 应用开发的“大爆炸”理论。如果我们不是迭代式地开发软件，而是在一场旷日持久的马拉松中写下大量代码，那我们就只能指望结束时的那次“大爆炸”能神奇地产生一个能工作的应用。遗憾的是，这种魔法很少发生。}

\section{变化与复杂度是良好设计的大敌}

变化与复杂度是良好设计的首要挑战，也是贯穿本书的两大主题。对软件开发人员来说，变化与复杂度是生活中不可避免的现实。任何没有把变化考虑在内的设计，都会在第一个变更请求到来时很快陷入麻烦。随着开发项目推进，代码可能变得凌乱而复杂，尤其是在项目团队中有多名程序员的时候。设计原则与设计模式正是用来应对变化与复杂度的工具。

变化可能发生在开发过程中。在我们编写代码时，应用需求可能会变化。或者我们可能改变了对设计的想法，想要重写部分代码。在应用完成并部署之后，我们可能会收到添加或修改功能的需求。良好的设计使我们能够修改代码的某一部分，而无需修改其他部分。良好的设计促进了应对变化所必需的灵活性。

然而，如果不加小心，试图让应用灵活到足以应对许多变化，反而可能增加它的复杂度。我们的程序可能变得复杂到无法管理。良好的设计还有助于防止应用失控。优秀的软件设计师必须经常做出权衡。

\section{用面向对象编程概念进行设计}

本书讨论的是面向应用开发的、被广泛使用的 OOP 范式的软件设计。OOP 概念是本书所涵盖的所有良好设计原则和设计模式的基础。因此，快速回顾一下这些概念，有助于确保我们大家达成共识。

OOP 基于以下四个主要概念：

\begin{itemize}
\item 封装
\item 抽象
\item 继承
\item 多态
\end{itemize}

\emph{封装}关乎类的使用。C++ 类通常包含成员变量和成员函数。在运行时，应用会从类中\emph{实例化}（创建）对象。在运行期间，对象会经历状态变化。对象的每一个状态都由其成员变量的一组唯一取值来表征。每当对象的成员变量取值发生变化，对象就会从一个状态转换到另一个状态。对象的成员函数可以作用于这些成员变量。这些函数决定了对象在运行时的\emph{行为}。

类的公有成员构成它的\emph{接口}，任何代码都可以访问；而它的私有成员只能由该类的成员函数访问。类的受保护成员只能由该类的子类的成员函数访问。

在本书中，我们会超越封装严格的 OOP 含义，用它来容纳并隔离应用中的变化。\emph{封装变化之处}是我们将多次使用的一条设计原则。第 5 章会讲到如何把类的成员设为私有和受保护，以封装对其实现的修改。

\myGraphic{0.755}{images/CH01_F03_UN07_Mak.png}{}

\emph{抽象}意味着忽略无关的细节，只关注与我们正在开发的应用相关的重点。恰当地使用抽象是降低复杂度的重要方式。

\emph{继承}使我们能够从超类（父类或基类）创建子类（派生类）。父类把状态信息（以成员变量的形式）和行为（以成员函数的形式）传递给它的子类。每个子类都可以添加自己额外的状态和行为，也可以重写（替换）任何继承来的状态或行为。

在运行时，如果你有一个指向某个对象的指针，而该对象是从一个共同超类的某个子类实例化而来的，那么\emph{多态}会根据是哪个子类实例化了该对象，来决定这个对象如何表现（即执行哪些成员函数）。多态有助于简化应用的设计。在整本书中，我们将看到 OOP 概念是本书所涵盖的设计原则和设计模式的基础。

\section*{小结}

\begin{itemize}
\item 设计是一种有纪律的工程方法，用于为某个问题创建解决方案。在软件工程中，问题在于创建可工作的软件，而解决方案就是一个设计良好、可持续的应用。
\item 设计良好的软件在很多方面都更胜一筹，例如更可靠、更灵活、更易维护。良好的设计有助于确保应用按时完成，并符合客户的预期。
\item 通过运用良好的软件设计技术（包括良好的设计原则和设计模式），是有可能成为更优秀的程序员的。
\item 良好的设计原则有助于让我们的代码更灵活，能够应对诸如新需求之类的变化。设计模式是经过业界验证的模型，用于为常见软件架构问题创建解决方案。
\item 软件设计始于获取并分析应用的需求，以确保我们开发的是正确的应用。应用必须完成它应该完成的工作。
\item 开发一个设计良好的可持续应用，几乎总是需要多轮迭代，并对糟糕的设计决策进行回溯。这是一项艰苦的工作。不要指望在一场马拉松式编码结束时会有神奇的“大爆炸”。
\item 良好的软件设计必须应对变化与复杂度这两大挑战。
\item 本书中的设计原则和设计模式基于面向对象编程的抽象、封装、继承和多态这几个概念。
\item 封装也意味着隔离程序中可能发生变化的部分。这样，当变化真的发生时，它们就不会泄漏出去，导致程序的其他部分也要跟着修改。
\end{itemize}
